%% ============================================================
%% Anhang
%% ============================================================

\chapter{Anhang}
\label{chap:anhang}

% Hinweis: Der Anhang enthaelt ergaenzendes Material, das im
% Haupttext den Lesefluss stoeren wuerde, aber fuer die
% Nachvollziehbarkeit der Arbeit wichtig ist.


\section{Interviewleitfaden}
\label{sec:interviewleitfaden}

% Falls Interviews durchgefuehrt wurden, hier den Leitfaden anfuegen.
% Beispiel:

% Das folgende Leitfadeninterview wurde am [Datum] mit [Person/Funktion]
% durchgefuehrt. Die Fragen gliederten sich in folgende Themenbloecke:
%
% \begin{enumerate}
%   \item Einstiegsfragen
%   \begin{itemize}
%     \item Koennten Sie sich kurz vorstellen?
%     \item Wie sind Sie zu diesem Thema gekommen?
%   \end{itemize}
%   \item Hauptfragen
%   \begin{itemize}
%     \item [Frage 1]
%     \item [Frage 2]
%   \end{itemize}
%   \item Abschlussfragen
%   \begin{itemize}
%     \item Gibt es etwas, das Sie noch ergaenzen moechten?
%   \end{itemize}
% \end{enumerate}


\section{Interviewtranskripte}
\label{sec:transkripte}

% Transkripte koennen hier eingefuegt oder als separate Dateien
% eingebunden werden:
% \input{kapitel/transkript-interview1}

[Hier die transkribierten Interviews einfuegen, falls vorhanden.]


\section{Erklaerung zum Einsatz von KI-Werkzeugen}
\label{sec:ki-erklaerung}

% Hinweis: Gemaess den aktuellen Richtlinien des BMBWF muss die
% Verwendung von KI-gestuetzten Werkzeugen (z.B. ChatGPT, DeepL,
% Grammarly) transparent dokumentiert werden.

Im Rahmen dieser Vorwissenschaftlichen Arbeit wurden folgende
KI-gestuetzte Werkzeuge eingesetzt:

\begin{itemize}
  \item \textbf{[Name des Werkzeugs, \zB ChatGPT 4o]:}
    [Beschreibung des Einsatzzwecks, \zB \glqq Wurde zur Ideenfindung
    und Strukturierung von Argumenten verwendet. Alle generierten
    Inhalte wurden kritisch geprueft und eigenstaendig ueberarbeitet.\grqq]
  \item \textbf{[Name des Werkzeugs, \zB DeepL]:}
    [Beschreibung des Einsatzzwecks, \zB \glqq Wurde zur Uebersetzung
    englischsprachiger Fachliteratur als Verstaendnishilfe genutzt.\grqq]
\end{itemize}

% Falls keine KI-Werkzeuge eingesetzt wurden:
% Im Rahmen dieser Vorwissenschaftlichen Arbeit wurden keine
% KI-gestuetzten Werkzeuge eingesetzt.


\section{Weitere Materialien}
\label{sec:weitere-materialien}

% Hier koennen weitere Materialien eingefuegt werden, z.B.:
% - Fragebogen
% - Statistische Auswertungen
% - Abbildungen in hoeherer Aufloesung
% - Quellenabdrucke

[Weitere ergaenzende Materialien hier einfuegen.]
